%---------------------------------------------------------------------------------------

\chapter[Revisão Bibliográfica]{Revisão Bibliográfica}\label{ch:revisao-bibliografica}

%---------------------------------------------------------------------------------------

% Lista de abreviaturas e de siglas
\criarsigla{DDF}{\textit{Double Distribution Function}}
\criarsigla{DNS}{\textit{Direct Numerical Simulation}}
\criarsigla{FHP}{Frisch-Hasslacher-Pomeau}
\criarsigla{FVM}{\textit{Finite Volume Method}}
\criarsigla{HPP}{Hardy-Pomeau-Pazzis}
\criarsigla{LES}{\textit{Large Eddy Simulation}}
\criarsigla{MPI}{\textit{Message Passing Interface}}
\criarsigla{PCIe}{\textit{Peripheral Component Interconnect Express}}
\criarsigla{TPU}{\textit{Tensor Processing Unit}}
\criarsigla{V\&V}{Verificação e Validação}

%---------------------------------------------------------------------------------------

% Lista de simbolos

%---------------------------------------------------------------------------------------

A simulação de alta fidelidade em dinâmica dos fluidos computacional (CFD) depende simultaneamente da adequação do modelo físico adotado, do rigor dos procedimentos que asseguram a confiabilidade dos resultados e da capacidade computacional disponível para obtê-los em prazo compatível com o uso pretendido \parencite{Moin2011, Oberkampf2010}.
Esses condicionantes raramente podem ser tratados de forma isolada, pois o refinamento exigido pela fidelidade eleva o custo do cálculo, e as arquiteturas capazes de absorver esse custo impõem restrições sobre os algoritmos que nelas se pretende executar \parencite{Karp2025, Kolla2018}.
O método de Lattice Boltzmann (LBM) ganhou espaço justamente por conciliar esses aspectos, ao substituir a discretização direta das equações de Navier-Stokes por uma dinâmica local de funções de distribuição que se adapta com naturalidade às máquinas de computação de alto desempenho (HPC) \parencite{Chen1998, Succi2001, Tomczak2021}.
Sua aplicação a escoamentos compressíveis, contudo, é relativamente recente e apoia-se em formulações que ainda coexistem na literatura \parencite{Tran2022, Rathnayake2024, Coreixas2017}, do mesmo modo que a escolha da tecnologia com que um resolvedor será paralelizado permanece uma decisão pouco amparada por evidências comparativas \parencite{Davis2023, Martineau2016}.

%---------------------------------------------------------------------------------------

\section{Simulações de Alta Fidelidade}\label{sec:simulacoes-de-alta-fidelidade}

Os fenômenos fluidodinâmicos de interesse prático apresentam forte dependência de pequenas escalas espaciais e temporais, o que exige discretizações refinadas e esquemas numéricos de alta ordem.
Ao mesmo tempo, ensaios experimentais em túnel de vento tornam-se custosos e, por vezes, inviáveis quando se busca explorar amplos espaços de parâmetros, como a variação do número de Mach, do número de Reynolds e da geometria do domínio.
O emprego de métodos numéricos que conciliem exatidão e eficiência computacional é, portanto, essencial em aplicações e pesquisas aeroespaciais e farmacêuticas, assim como no desenvolvimento de turbomáquinas e de motores de combustão \parencite{Mira2023, Succi2001}.
A versatilidade do LBM no tratamento de escoamentos multifásicos e de meios porosos tem sido amplamente documentada, consolidando o método como alternativa competitiva aos resolvedores tradicionais baseados no método de volumes finitos (\textit{Finite Volume Method}, FVM) \parencite{Lauricella2025, Aniello2022}.

Os métodos tradicionais de CFD de alta fidelidade, como a simulação de grandes escalas (\textit{Large Eddy Simulation}, LES) e a simulação numérica direta das equações de Navier-Stokes (\textit{Direct Numerical Simulation}, DNS), alcançam elevada precisão, porém são tipicamente limitados pelo custo computacional e pelo consumo de memória.
Isso ocorre porque a representação dos efeitos turbulentos associados à vorticidade do escoamento requer resolver um amplo intervalo de escalas, desde a escala característica do problema até a escala de Kolmogorov, na qual os vórtices são dissipados por efeitos viscosos \parencite{Davidson2015}.
São necessárias, assim, malhas com dezenas de milhões de elementos, o que impõe grande tráfego de memória e elevada capacidade de processamento para preservar a estabilidade numérica e a captura das pequenas escalas.
Por essa razão, o DNS é viável apenas para números de Reynolds reduzidos, enquanto o LES emprega modelos de filtragem para as pequenas escalas, sendo mais frequentemente adotado quando simulações de alta fidelidade são requeridas \parencite{Karp2025}.

Para contornar esses desafios, técnicas de redução de custo e de ampliação do paralelismo são imprescindíveis.
A divisão do domínio em subdomínios tratados em nós distintos de processamento, associada a esquemas de decomposição de domínio e a algoritmos eficientes de comunicação, permite executar simulações de grande porte com alta resolução em \textit{hardware} moderno, reduzindo o tempo de processamento e o custo computacional \parencite{Kolla2018}.
\textcite{Karp2025} demonstra que, mesmo em arquiteturas modulares que combinam CPUs e GPUs, é possível manter eficiência paralela superior a 80\% em milhares de núcleos, desde que a decomposição do domínio e a comunicação entre processos sejam cuidadosamente projetadas.

%---------------------------------------------------------------------------------------

\subsection{Análise de Erros, Verificação e Validação}\label{subsec:erros-verificacao-validacao}

A confiabilidade das simulações de alta fidelidade em CFD fundamenta-se em processos rigorosos de verificação e validação (V\&V) \parencite{Roache1998, Oberkampf2010}.
Enquanto a verificação busca assegurar que o modelo matemático está sendo resolvido corretamente, ou seja, que as equações são resolvidas de forma correta, a validação visa confirmar se o modelo escolhido representa fielmente a realidade física, isto é, se as equações resolvidas são as adequadas ao fenômeno \parencite{AIAA1998, ASME2009}.

A análise de fenômenos fluidodinâmicos pode ser conduzida por via experimental ou por simulação numérica, cada qual com vantagens e limitações distintas.
A abordagem experimental, usualmente realizada em túneis de vento ou em bancadas de ensaio, oferece contato direto com a realidade física, mas enfrenta o alto custo de operação, a dificuldade de acesso a dados em todo o domínio e a presença de erros de medição, de calibração e de interferência dos sensores \parencite{Stern2001}.
A simulação numérica, em contrapartida, permite extrair dados detalhados em qualquer ponto do domínio e explorar amplos espaços de parâmetros a custo comparativamente reduzido, ainda que sua precisão dependa intrinsecamente da qualidade dos modelos matemáticos e da resolução numérica empregada \parencite{Oberkampf2010}.

O processo de simulação inicia-se com a modelagem matemática, na qual o problema físico é traduzido em equações diferenciais governantes, como as equações de Navier-Stokes ou a equação de Boltzmann.
Simplificações são frequentemente introduzidas para tornar o problema tratável, entre elas as hipóteses de fluido incompressível, de gás ideal ou de regime estacionário.
A solução dessas equações pode ser obtida analiticamente, fornecendo resultados exatos, porém restritos a geometrias e condições de contorno elementares.
Para escoamentos complexos de interesse prático, a solução numérica torna-se imperativa, aproximando-se as equações contínuas por meio de métodos de discretização.

A transição entre a realidade física e o resultado numérico introduz diversas fontes de erro que afetam a fidelidade da análise.
O erro de modelagem decorre das simplificações nas equações governantes ou do uso de modelos que não capturam toda a física envolvida, manifestando-se na incapacidade de reproduzir observações experimentais \parencite{AIAA1998}.
O erro de discretização surge da transformação das equações diferenciais em algébricas, englobando truncamentos espaciais e temporais que afetam a convergência e a captura das pequenas escalas \parencite{Roache1998}.
Somam-se a esses o erro de iteração, associado ao processo de solução numérica até que os critérios de convergência sejam atingidos, e o erro de arredondamento, decorrente da precisão finita da representação numérica.
No contexto do LBM, a verificação exige estudos de convergência de malha que demonstrem a redução do erro global com taxa de segunda ordem em escoamentos subsônicos \parencite{Rathnayake2024, Succi2001}.
Em regimes compressíveis e hipersônicos, a presença de choques e descontinuidades introduz erros adicionais capazes de degradar a ordem de convergência, o que exige técnicas de estabilização e sensores de choque para preservar a integridade dos resultados \parencite{Tran2022, Qiu2020, Guo2020}.

%---------------------------------------------------------------------------------------

\section{Método de Lattice Boltzmann}\label{sec:metodo-de-lattice-boltzmann}

O método de Lattice Boltzmann consolidou-se nas últimas décadas como abordagem mesoscópica robusta para a simulação de escoamentos, oferecendo alternativa computacionalmente eficiente aos métodos baseados na discretização direta das equações de Navier-Stokes.
Além de simplificar o tratamento de geometrias complexas e de condições de contorno, o LBM integra-se naturalmente a arquiteturas de paralelismo massivo em razão do caráter local de suas operações \parencite{Chen1998, Succi2001, WolfGladrow2000}.
A fundamentação do método repousa na teoria cinética dos gases, modelando a evolução de funções de distribuição de partículas sobre uma rede discreta em vez de resolver diretamente as variáveis macroscópicas contínuas.

%---------------------------------------------------------------------------------------

\subsection{Evolução Histórica e Autômatos de Gás de Rede}\label{subsec:lbm-historico}

A origem do LBM está associada aos autômatos de gás de rede (LGA).
O modelo precursor, conhecido como HPP (Hardy, Pomeau e Pazzis), demonstrou que leis macroscópicas de conservação podem emergir de regras locais simples de colisão e propagação em uma rede bidimensional quadrada \parencite{HPP1973}.
Esse modelo, no entanto, carecia de isotropia rotacional, o que impedia a recuperação correta das equações de Navier-Stokes.
O obstáculo foi superado por \textcite{Frisch1986} com a introdução do modelo FHP (Frisch, Hasslacher e Pomeau), que empregava uma rede hexagonal para assegurar a isotropia necessária, marcando o início da dinâmica dos fluidos computacional baseada em redes discretas.

Como discutido por \textcite{Chen1998}, embora os modelos LGA fossem teoricamente atraentes por sua simplicidade e por sua conservação estrita e incondicional, apresentavam desafios práticos significativos.
Entre os principais obstáculos estavam o ruído estatístico inerente às variáveis booleanas que representavam a presença ou a ausência de partículas, a falta de invariância de Galileu e uma dependência anômala da pressão com a velocidade.
A transição para o LBM ocorreu quando se propôs substituir as partículas booleanas discretas pelas médias de conjunto dessas variáveis, $f_i$, funções de distribuição contínuas que representam a probabilidade de encontrar uma partícula em determinada direção de velocidade.

Esse avanço, liderado por \textcite{McNamara1988}, eliminou o ruído estatístico associado aos métodos LGA e permitiu representação mais fiel da hidrodinâmica.
Posteriormente, \textcite{Higuera1989} introduziram a linearização do operador de colisão em torno da distribuição de equilíbrio, o que facilitou a análise de estabilidade e viabilizou a implementação numérica que fundamenta a maioria das variantes atuais do método.
Em seguida, modelos de tempo de relaxação único (BGK) foram propostos de forma independente e consolidados na literatura por \textcite{Qian1992}, fornecendo a formulação simples e amplamente utilizada do LBM contemporâneo.

%---------------------------------------------------------------------------------------

\subsection{Escalas de Descrição e Modelagem Mesoscópica}\label{subsec:lbm-escalas}

A descrição de sistemas fluidodinâmicos pode ser abordada sob três perspectivas, conforme a escala de observação e o grau de detalhamento requerido.
Na escala molecular, ou microscópica, o fluido é tratado como um conjunto de partículas individuais que seguem as leis do movimento de Newton, o que exige o rastreamento de trajetórias individuais por meio de simulações de dinâmica molecular.
Ainda que precisa, essa descrição é computacionalmente proibitiva para os sistemas macroscópicos associados a problemas clássicos de engenharia.
No extremo oposto, a escala macroscópica trata o fluido como meio contínuo de propriedades médias, regido pelas equações de conservação e pelas equações de Navier-Stokes.

O método de Lattice Boltzmann situa-se na escala mesoscópica, servindo como ponte entre as descrições microscópica e macroscópica.
Em vez de rastrear cada partícula, o método opera sobre a função de distribuição de probabilidade $f(\mathbf{x}, \boldsymbol{\xi}, t)$, que descreve a densidade provável de partículas localizadas na posição $\mathbf{x}$ com velocidade microscópica $\boldsymbol{\xi}$ no instante $t$.
A evolução dessa função é governada pela equação de transporte de Boltzmann:
\begin{equation}
    \frac{\partial f}{\partial t} + \boldsymbol{\xi} \cdot \nabla f + \frac{\mathbf{F}}{\rho} \cdot \nabla_{\boldsymbol{\xi}} f = \Omega(f)
\end{equation}
onde $\mathbf{F}$ representa as forças externas volumétricas que atuam sobre o fluido, como a gravidade, $\nabla_{\boldsymbol{\xi}}$ é o gradiente no espaço de velocidades e $\Omega(f)$ denota o operador de colisão, que descreve as variações de $f$ decorrentes dos choques entre partículas.

A \autoref{fig:lbm-escalas} ilustra a hierarquia dessas escalas e o posicionamento do LBM como elo entre a descrição molecular e o comportamento contínuo.

\begin{figure}[H]
    \centering
    \caption{Escalas de descrição de um fluido: microscópica, com partículas individuais e dinâmica molecular; mesoscópica, com funções de distribuição de probabilidade e equação de Boltzmann; e macroscópica, com meio contínuo e equações de Navier-Stokes.}
    \label{fig:lbm-escalas}
    \fbox{\begin{minipage}{0.85\textwidth}
        \centering
        \vspace{2cm}
        [Espaço reservado para figura: hierarquia das escalas de descrição de um fluido, com três blocos dispostos da esquerda para a direita. No primeiro, a escala microscópica, com partículas individuais e trajetórias governadas pelas leis de Newton. No segundo, a escala mesoscópica, com a função de distribuição de probabilidade sobre uma rede de velocidades discretas e a equação de Boltzmann. No terceiro, a escala macroscópica, com o fluido representado como meio contínuo e as equações de Navier-Stokes. Setas entre os blocos devem indicar as transições, associadas à promediação estatística e à expansão de Chapman-Enskog.]
        \vspace{2cm}
    \end{minipage}}
    \fonte{Adaptado de \textcite{Chen1998}.}
\end{figure}

A transição quantitativa entre essas escalas é parametrizada pelo número de Knudsen, $Kn = \lambda/L$, em que $\lambda$ é o percurso livre médio das partículas e $L$ é o comprimento característico do escoamento.
O LBM é aplicável no limite de meio contínuo ($Kn \to 0$), no qual recupera formalmente as equações de Navier-Stokes por meio da análise assintótica multiescala de Chapman-Enskog \parencite{Chen1998, Benzi1992}.
A análise multiescala que sustenta essa recuperação foi originalmente desenvolvida por \textcite{Chapman1916} e por \textcite{Enskog1917} para o transporte em gases rarefeitos, e sua sistematização de referência encontra-se em \textcite{Chapman1970}.
Sua adaptação à equação de Boltzmann previamente discretizada no espaço de velocidades é o que fornece as relações entre o tempo de relaxação e os coeficientes de transporte empregadas no LBM.
A demonstração que conecta a equação de Boltzmann contínua à sua forma discretizada no espaço e no tempo foi formalizada por \textcite{He1997}.

%---------------------------------------------------------------------------------------

\subsection{Pesquisas em Escoamentos Compressíveis e Ondas de Choque}\label{subsec:lbm-pesquisas-compressivel}

A aplicação do LBM a escoamentos compressíveis acompanha o método desde seus primeiros anos, mas somente nas duas últimas décadas alcançou maturidade suficiente para tratar choques intensos.
A restrição de origem é conhecida e decorre de duas causas.
O equilíbrio obtido pelo truncamento da distribuição de Maxwell-Boltzmann em segunda ordem na velocidade macroscópica limita o método a números de Mach reduzidos, e a rede de velocidades discretas usual não possui momentos isotrópicos de ordem suficientemente elevada para reproduzir o fluxo de energia do regime térmico.
As primeiras tentativas de superá-la introduziram uma segunda população para a energia interna, como no modelo termo-hidrodinâmico de \textcite{Alexander1993}, que já obtinha variação de temperatura acoplada ao campo de velocidade, ainda que restrita a pequenas perturbações.
A revisão de \textcite{Chen1998} registra esse estado da arte no fim da década de 1990, apontando a compressibilidade artificial associada ao erro de ordem $Ma^2$ e a ausência de um modelo térmico estável como as principais limitações então em aberto.

Ainda naquele período, duas linhas de solução se delinearam.
A primeira consistiu em ampliar o conjunto de velocidades discretas, dando origem às redes multivelocidade, nas quais se torna possível satisfazer as condições de isotropia dos momentos de ordem superior exigidas pelo regime térmico.
Nessa direção, \textcite{Sun1998} propôs um conjunto de velocidades adaptativo, ajustado localmente à velocidade e à temperatura do escoamento, e \textcite{Yan1999} construíram um modelo capaz de capturar ondas de choque com um número reduzido de velocidades.
A segunda substituiu a propagação exata sobre a rede por esquemas de diferenças finitas aplicados à equação de Boltzmann discreta, estratégia adotada por \textcite{Watari2003} para obter elevada isotropia espacial em um modelo térmico bidimensional e por \textcite{Kataoka2004} para recuperar as equações de Euler compressíveis.
Na mesma direção, \textcite{Qu2007} abandonaram a forma polinomial do equilíbrio e passaram a construí-lo por integração de uma função circular sobre o espaço de velocidades, o que estendeu a validade do modelo a escoamentos invíscidos com números de Mach elevados.

A partir de meados da década de 2000, o esforço deslocou-se para a estabilização do método sobre redes de tamanho convencional, mais convenientes do ponto de vista computacional.
\textcite{Latt2006} propuseram a regularização das funções de distribuição antes da colisão, substituindo a parcela de não equilíbrio por sua reconstrução a partir de um conjunto reduzido de momentos, procedimento que remove as contribuições espúrias de ordem superior responsáveis pela perda de estabilidade.
A ideia foi sistematizada por \textcite{Coreixas2017} na forma de uma regularização recursiva, em que os momentos de ordem superior são obtidos por relações de recorrência a partir dos momentos conservados, e por \textcite{Feng2019} em um modelo térmico híbrido, no qual a equação da energia é resolvida por diferenças finitas enquanto massa e quantidade de movimento permanecem sob a dinâmica de Lattice Boltzmann.
Em paralelo, a formulação entrópica foi estendida ao regime compressível por \textcite{Frapolli2015}, que determinam o equilíbrio pela minimização de um funcional de entropia sujeito às restrições de conservação, e \textcite{Saadat2019} mostraram ser possível ajustar de forma independente o número de Prandtl e o expoente adiabático mantendo a rede padrão.

O tratamento de escoamentos com choques fortes e grandes variações de temperatura consolidou-se no período de 2018 a 2025 em torno de duas ideias complementares.
A primeira é a adoção de referenciais móveis para o conjunto de velocidades discretas, proposta por \textcite{Dorschner2018} sob a designação de partículas sob demanda, na qual cada nó da malha define seu próprio referencial, com extensão posterior a escoamentos com descontinuidades severas \parencite{Kallikounis2022}.
A segunda é o emprego de duas funções de distribuição (\textit{Double Distribution Function}, DDF), uma para massa e quantidade de movimento e outra para a energia total, associado a redes deslocadas (\textit{shifted lattices}) cujo deslocamento acompanha a velocidade local do escoamento.
Essa combinação foi formulada por \textcite{Tran2022} para escoamentos compressíveis com elevado número de Reynolds e posteriormente detalhada e validada por \textcite{Rathnayake2024} em casos supersônicos, com conservação de energia e estabilidade termodinâmica preservadas na presença de choques.
Entre as demais contribuições recentes situam-se a formulação semilagrangiana de \textcite{Wilde2020}, que remove a restrição de que o deslocamento das populações coincida com os nós da malha, e os esquemas híbridos aperfeiçoados de \textcite{Renard2021}, aplicados a regimes subsônicos e supersônicos sobre redes padrão.
A validação dessas técnicas para a captura de ondas de choque em escoamentos transônicos e supersônicos foi ampliada por \textcite{Qiu2020}, e \textcite{Guo2020} demonstraram sua aplicação a escoamentos tridimensionais com grandes variações nos campos de temperatura e de massa específica.

O panorama resultante é o de um campo em que nenhuma formulação se tornou dominante, e no qual a escolha depende do compromisso aceito entre custo por nó da malha, faixa de números de Mach coberta e complexidade de implementação.
Esses avanços, em conjunto, consolidam o LBM como alternativa viável para simulações aeroespaciais complexas, ao mesmo tempo que elevam de forma expressiva o custo do ciclo de cálculo em relação ao método isotérmico clássico, o que reforça a importância da eficiência da implementação computacional.

%---------------------------------------------------------------------------------------

\section{Tecnologias de Computação de Alto Desempenho}\label{sec:tecnologias-de-computacao-de-alto-desempenho}

A evolução dos supercomputadores em direção a arquiteturas heterogêneas favoreceu particularmente algoritmos como o LBM, dada a localidade espacial da etapa de colisão e o padrão regular de acesso à memória da etapa de propagação \parencite{Ma2023, Tomczak2021}.
As GPUs, originalmente desenvolvidas para renderização gráfica, mostraram-se eficientes para o cálculo científico em razão de sua capacidade de executar milhares de operações simultâneas, enquanto as FPGAs oferecem flexibilidade para a implementação de algoritmos específicos com elevada eficiência energética \parencite{Karp2025, Herdman2012}.
Somam-se a esses dispositivos os aceleradores dedicados a operações tensoriais (\textit{Tensor Processing Unit}, TPU), que ampliam ainda mais a diversidade de alvos de execução disponíveis.
A \autoref{fig:hpc-arquitetura-heterogenea} resume a organização típica de um nó de computação heterogêneo e os respectivos níveis de paralelismo explorados pelos modelos de programação discutidos nesta seção.

\begin{figure}[H]
    \centering
    \caption{Organização de um nó de computação heterogêneo e níveis de paralelismo explorados pelos modelos de programação avaliados neste trabalho.}
    \label{fig:hpc-arquitetura-heterogenea}
    \fbox{\begin{minipage}{0.85\textwidth}
        \centering
        \vspace{2cm}
        [Espaço reservado para figura: esquema de um nó de computação heterogêneo, com uma CPU multinúcleo e sua memória principal conectadas por um barramento PCIe a um ou mais aceleradores com memória dedicada. O esquema deve indicar, ao lado de cada elemento, o modelo de programação associado: OpenMP e MPI na CPU, CUDA, OpenCL, OpenACC, Kokkos e RAJA nos aceleradores, e MPI na comunicação entre nós de um agrupamento.]
        \vspace{2cm}
    \end{minipage}}
    \fonte{O autor (2026).}
\end{figure}

A crescente diversidade de arquiteturas levou ao desenvolvimento de tecnologias de paralelização que auxiliam o programador a escrever códigos executáveis de forma eficiente em plataformas heterogêneas.
A escolha entre modelos proprietários, como o CUDA, e camadas de abstração, como o Kokkos e o RAJA, envolve um compromisso entre desempenho e portabilidade.
Os modelos proprietários tendem a oferecer otimizações de baixo nível que exploram ao máximo um \textit{hardware} específico, resultando em desempenho superior para aplicações críticas, mas restringem a portabilidade do código e exigem reescritas significativas na migração para arquiteturas distintas.
As camadas de abstração, por sua vez, permitem que o mesmo código seja executado em múltiplas plataformas com modificações mínimas, ainda que essa portabilidade possa vir acompanhada de leve penalização de desempenho \parencite{Wong2016}.

Essa comparação é determinante para a escolha dos modelos de programação empregados no desenvolvimento dos miniaplicativos apresentados neste trabalho, cuja meta é avaliar não apenas o desempenho computacional, mas também o custo de desenvolvimento e a escalabilidade do código em arquiteturas heterogêneas.
Os principais modelos de programação utilizados em HPC são examinados a seguir, agrupados segundo a filosofia com que expõem o paralelismo ao programador, com ênfase em suas origens, nas arquiteturas suportadas e em suas características distintivas.
Cada um deles é retomado no Capítulo~\ref{ch:implementacao-computacional} do ponto de vista de sua implementação.

%---------------------------------------------------------------------------------------

\subsection{Modelos Baseados em Diretivas de Compilação}\label{subsec:modelos-por-diretivas}

Os modelos baseados em diretivas partem do pressuposto de que o código sequencial já existe e de que o paralelismo pode ser expresso por anotações interpretadas pelo compilador.
A principal consequência dessa filosofia é a preservação de uma única base de código: quando as diretivas são ignoradas, o programa permanece válido e produz o mesmo resultado em execução sequencial.

Concebido para facilitar a programação em sistemas de memória compartilhada, o OpenMP (\textit{Open Multi-Processing}) é amplamente utilizado em CPUs multinúcleo.
Seu modelo de execução é o de bifurcação e junção (\textit{fork-join}), no qual um fluxo de execução principal cria um conjunto de fluxos auxiliares ao entrar em uma região paralela e os reúne ao final dela.
A partir da versão 4.0, a especificação passou a oferecer suporte também a dispositivos aceleradores por meio das diretivas de descarregamento (\textit{offloading}), o que ampliou seu alcance para além da memória compartilhada.
Seus principais atrativos, em relação a outros modelos, são a facilidade de integração, a maturidade da especificação e a disponibilidade em praticamente todos os compiladores de uso corrente \parencite{Dagum1998}.

O OpenACC (\textit{Open Accelerators}) é um modelo de programação também baseado em diretivas, desenvolvido para simplificar a aceleração de trechos de código em dispositivos heterogêneos, especialmente GPUs.
Ao permitir que o programador indique quais regiões do código devem ser paralelizadas, o OpenACC facilita a adaptação de aplicações existentes a plataformas de HPC sem a reescrita completa do código.
Sua abordagem é descritiva, e não prescritiva: o programador informa a existência de paralelismo e delega ao compilador a decisão sobre como mapear as iterações no dispositivo, o que reduz o esforço de programação, mas torna o desempenho dependente da qualidade do compilador empregado.
Sua especificação é mantida por um consórcio de empresas e instituições acadêmicas, e sua eficácia tem sido avaliada em diversos estudos comparativos \parencite{Wolfe2018}.

%---------------------------------------------------------------------------------------

\subsection{Modelos de Baixo Nível para Aceleradores}\label{subsec:modelos-de-baixo-nivel}

Os modelos de baixo nível adotam a filosofia oposta à dos modelos por diretivas.
O paralelismo passa a ser expresso explicitamente por meio de funções de cálculo executadas no dispositivo, denominadas \textit{kernels}, e o programador assume o controle da hierarquia de memória, da organização dos fluxos de execução e da sincronização entre o sistema principal e o acelerador.

A plataforma CUDA, desenvolvida pela NVIDIA, representa um dos modelos proprietários mais maduros para programação em GPUs.
Sua arquitetura permite controle fino sobre a memória e sobre a execução paralela, resultando em desempenho elevado para aplicações otimizadas.
O modelo de execução organiza os fluxos em blocos e grades, e disponibiliza níveis distintos de memória, entre eles a memória global, a memória compartilhada por bloco, a memória constante e os registradores, cuja utilização criteriosa é determinante para o desempenho final.
A natureza proprietária do CUDA, contudo, impõe desafios de portabilidade, uma vez que o código produzido restringe-se a \textit{hardware} NVIDIA \parencite{Chakrabarti2012}.

Mantido pelo Khronos Group, o OpenCL (\textit{Open Computing Language}) é uma interface de programação de aplicações aberta que possibilita o cálculo paralelo em diferentes dispositivos, incluindo CPUs, GPUs, FPGAs e outros aceleradores.
Sua arquitetura, baseada em \textit{kernels} e filas de comandos, permite escrever código portável, embora otimizações específicas para cada \textit{hardware} sejam frequentemente necessárias \parencite{Stone2010}.
Uma particularidade do padrão é a compilação dos \textit{kernels} em tempo de execução, o que permite adaptar o código ao dispositivo efetivamente encontrado na máquina, ao custo de uma latência adicional na inicialização do programa e da manutenção dos \textit{kernels} em arquivos separados do restante da aplicação.

%---------------------------------------------------------------------------------------

\subsection{Camadas de Abstração para Portabilidade de Desempenho}\label{subsec:camadas-de-abstracao}

As camadas de abstração surgiram como resposta à fragmentação do ecossistema de aceleradores nos grandes laboratórios de pesquisa, cujos códigos precisam sobreviver à substituição periódica das máquinas e, portanto, à troca de fabricante.
A estratégia comum a essas bibliotecas consiste em expressar o paralelismo por meio de construtores de C++ resolvidos em tempo de compilação, de modo que a mesma fonte gere código nativo para o alvo selecionado sem penalidade de indireção em tempo de execução.

Desenvolvido no Sandia National Laboratories, o Kokkos é uma camada de abstração para programação paralela que visa à portabilidade entre arquiteturas, incluindo CPUs e GPUs, sem prejuízo de desempenho.
O Kokkos fornece abstrações que permitem escrever código em C++ de forma agnóstica à plataforma, delegando à biblioteca a otimização para o \textit{hardware} disponível \parencite{CarterEdwards2014}.
Seu diferencial está em abstrair simultaneamente a execução, por meio de espaços de execução, e os dados, por meio de estruturas multidimensionais cuja disposição em memória é escolhida automaticamente conforme o alvo, favorecendo a vetorização em CPUs e o acesso coalescido em GPUs.

O RAJA, desenvolvido no Lawrence Livermore National Laboratory, é outra camada de abstração voltada à escrita de código C++ portável para arquiteturas heterogêneas.
De forma semelhante ao Kokkos, o RAJA abstrai as complexidades da programação paralela, permitindo que o desenvolvedor concentre-se na lógica do algoritmo em vez dos detalhes específicos do \textit{hardware} \parencite{Beckingsale2019}.
Sua abrangência é, entretanto, deliberadamente menor: o RAJA abstrai apenas os laços de repetição, por meio de políticas de execução, e mantém sob responsabilidade da aplicação a decisão sobre onde os dados residem, o que facilita sua adoção incremental em códigos já existentes com estruturas de dados próprias.
Ambas as bibliotecas compõem-se com o padrão SYCL e com as extensões de fabricantes concorrentes, o que amplia o conjunto de aceleradores alcançáveis a partir de uma única base de código.

%---------------------------------------------------------------------------------------

\subsection{Paralelismo em Memória Distribuída}\label{subsec:paralelismo-distribuido}

Os modelos discutidos até aqui atuam no interior de um único nó de computação, no qual todos os fluxos de execução compartilham o mesmo espaço de endereçamento.
Simulações de grande porte, contudo, excedem a memória e a capacidade de processamento de um único nó, o que exige distribuir o domínio entre máquinas que se comunicam por uma rede de interconexão.

A interface de passagem de mensagens (\textit{Message Passing Interface}, MPI) é o padrão para comunicação entre processos em sistemas distribuídos.
Amplamente utilizada em agrupamentos de computadores (\textit{clusters}) e em supercomputadores, a MPI permite coordenar o cálculo entre múltiplos nós de processamento, viabilizando simulações em escalas massivas.
Cada processo detém um subdomínio e troca com seus vizinhos apenas as informações das camadas de fronteira, o que torna o custo de comunicação proporcional à área das interfaces e não ao volume do subdomínio, propriedade que sustenta a escalabilidade do modelo em problemas de grande porte.
Sua robustez e escalabilidade são comprovadas em inúmeros estudos e implementações, e sua especificação é gerida pelo MPI Forum, consórcio internacional de especialistas \parencite{TheMPIForum1993}.
Na prática corrente dos grandes centros de computação, a MPI não compete com os modelos anteriores, mas combina-se a eles em uma abordagem híbrida, na qual a MPI coordena os nós e um modelo de memória compartilhada ou de acelerador explora o paralelismo no interior de cada um deles \parencite{Kolla2018}.

%---------------------------------------------------------------------------------------

\subsection{Miniaplicativos como Instrumento de Avaliação}\label{subsec:miniaplicativos}

A avaliação comparativa de tecnologias de paralelização enfrenta uma dificuldade metodológica evidente: programas de aplicação completos acumulam, ao longo de anos de desenvolvimento, otimizações específicas de uma tecnologia, o que impede atribuir as diferenças observadas exclusivamente ao modelo de programação.
Reescrever integralmente um resolvedor maduro em cada tecnologia avaliada é, por sua vez, inviável no prazo de uma pesquisa.
Esse impasse pode ser resolvido com o desenvolvimento de miniaplicativos, programas deliberadamente reduzidos que preservam o núcleo de cálculo, o padrão de acesso à memória e o padrão de comunicação da aplicação de origem, mas descartam as funcionalidades acessórias.

O interesse dessa abordagem está em sua dupla função.
Como instrumento de medida, o miniaplicativo isola o comportamento que se pretende observar e permite reimplementá-lo integralmente em cada tecnologia com esforço compatível com um estudo comparativo, conforme demonstrado por \textcite{Herdman2012} na avaliação de modelos por diretivas frente a modelos de baixo nível.
Como instrumento de projeto, o miniaplicativo antecipa decisões de arquitetura de \textit{software} antes que elas sejam onerosas, funcionando como protótipo do resolvedor que se pretende construir.
Estudos recentes de portabilidade de desempenho apoiam-se justamente nessa estratégia para comparar camadas de abstração em máquinas de arquiteturas distintas \parencite{Davis2023, Karp2025}.

A prática consolidou-se a partir do conjunto Mantevo, organizado no Sandia National Laboratories, que reuniu miniaplicativos representativos de classes distintas de cálculo científico, entre eles o miniFE, para sistemas lineares esparsos oriundos de elementos finitos, o miniMD, para dinâmica molecular, e o miniGhost, para a troca de camadas de fronteira em malhas estruturadas \parencite{Heroux2009}.
No domínio da dinâmica dos fluidos e da hidrodinâmica computacional, três exemplos aproximam-se particularmente do presente trabalho.
O LULESH, mantido no Lawrence Livermore National Laboratory, reproduz o núcleo de um resolvedor lagrangiano de hidrodinâmica em malha não estruturada e tornou-se referência em estudos de portabilidade por dispor de implementações em praticamente todos os modelos de programação em uso corrente \parencite{Karlin2013}.
O CloverLeaf, desenvolvido no âmbito dos programas britânicos de computação de alto desempenho, resolve as equações de Euler compressíveis bidimensionais em malha cartesiana por um esquema explícito de volumes finitos, o que o aproxima do escopo físico desta pesquisa, e foi justamente sobre ele que \textcite{Herdman2012} confrontaram OpenACC, OpenCL e CUDA \parencite{Mallinson2013}.
O TeaLeaf, por sua vez, isola a solução implícita de um problema de condução de calor e permite comparar resolvedores iterativos e suas exigências de comunicação coletiva \parencite{McIntoshSmith2017}.

Dois casos particulares delimitam os extremos do espectro de complexidade dos miniaplicativos.
Em um deles, o BabelStream reduz-se a um conjunto de operações elementares sobre vetores, com o propósito único de medir a largura de banda de memória efetivamente alcançada por cada modelo de programação em cada arquitetura, informação de referência incontornável para algoritmos limitados por memória, como o LBM \parencite{Deakin2018}.
No outro, o HPCG substituiu a fatoração densa como métrica de classificação de supercomputadores precisamente porque seu padrão de acesso irregular e sua dependência de comunicação representam melhor as aplicações científicas reais \parencite{Dongarra2016}.
Entre esses extremos situam-se os núcleos de cálculo extraídos de resolvedores completos, como os empregados nas avaliações de desempenho do LBM em GPUs \parencite{Tomczak2021, Ma2023}, e os protótipos que precederam resolvedores de produção em métodos de elementos espectrais \parencite{Jansson2024}.
O miniaplicativo desenvolvido nesta pesquisa segue essa mesma linha, com a particularidade de reproduzir integralmente a formulação compressível de interesse, e não apenas seu núcleo isotérmico.

A principal limitação da abordagem decorre de sua própria simplificação.
Por não reproduzir a diversidade de padrões de acesso de uma aplicação completa, o miniaplicativo pode subestimar efeitos de pressão sobre a hierarquia de memória ou de desequilíbrio de carga que só se manifestam em problemas reais, de modo que suas conclusões devem ser interpretadas como indicativas de tendências, e não como previsões absolutas de desempenho.

%---------------------------------------------------------------------------------------

\section{Síntese da Revisão e Lacuna Identificada}\label{sec:sintese-e-lacuna}

A literatura revisada permite três constatações que, em conjunto, delimitam a contribuição pretendida por este trabalho.
A primeira é que o LBM alcançou maturidade suficiente para tratar escoamentos compressíveis com ondas de choque, desde que associado a operadores de colisão robustos, a técnicas de regularização e a formulações de dupla distribuição \parencite{Tran2022, Rathnayake2024, Guo2020}.
A segunda é que o caráter local da colisão e a regularidade do acesso à memória na propagação tornam o método particularmente adequado às arquiteturas heterogêneas atuais, com desempenho reportado que se aproxima do teto imposto pela largura de banda de memória \parencite{Tomczak2021, Ma2023}.
A terceira é que existe hoje um número considerável de tecnologias de paralelização concorrentes, cada uma com um compromisso próprio entre desempenho, portabilidade e esforço de desenvolvimento.

Os estudos comparativos disponíveis, entretanto, concentram-se em núcleos de cálculo simples ou em versões isotérmicas do LBM, e raramente confrontam, sob um mesmo algoritmo e um mesmo caso de teste, o conjunto completo formado por modelos por diretivas, modelos de baixo nível, camadas de abstração e paralelismo distribuído.
Menos frequente ainda é a associação entre as métricas de desempenho e uma medida objetiva do esforço de implementação, informação essencial para quem precisa decidir qual tecnologia adotar no início de um projeto.
Essa lacuna justifica a metodologia adotada no Capítulo~\ref{ch:metodologia}, na qual a mesma formulação do LBM compressível é implementada em sete tecnologias distintas e avaliada simultaneamente quanto ao desempenho, à escalabilidade e à complexidade do código produzido.
